\clearpage\begingroup\NoSourceHead\SourceLeaf{511}
\renewcommand{\thefootnote}{\arabic{footnote})}\setcounter{footnote}{2}
\ArithmeticHead{54}{L}{BOLZA, GAUSS UND DIE VARIATIONSRECHNUNG.}
\begin{center}\textls[100]{b) Die Fläche ist in Parameterdarstellung gegeben.}\end{center}
\noindent\hspace*{\parindent}4. Bei der zweiten Herleitung der Differentialgleichung der geodätischen
Linien, die \textsc{Gauss} in Art. 19 gibt, wird die Fläche in Parameterdarstellung vorausgesetzt,
und es handelt sich um ein Variationsproblem vom einfachsten
Typus ohne Nebenbedingungen
\begin{equation*}\tag{4}
\int\sqrt{E\,dp^2+2F\,dp\,dq+G\,dq^2}=\text{Minimum.}
\end{equation*}

\textsc{Gauss} wählt $q$ als unabhängige Variable und erhält so die folgende Differentialgleichung
für $p$ als Funktion von $q$
\begin{equation*}\tag{5}
\frac{\partial E}{\partial p}\,dp^2+2\frac{\partial F}{\partial p}\,dp\,dq+\frac{\partial G}{\partial p}\,dq^2
=2\,ds\mathpunct{.}d\frac{E\,dp+F\,dq}{ds},
\end{equation*}
die er nach Ausführung der Differentiation und Einführung des Winkels $\theta$,
den die geodätische Linie mit der $p$-Richtung bildet, auf die Form bringt
\begin{equation*}\tag{6}
\sqrt{EG-F^2}\,d\theta=\tfrac12\frac{F}{E}\,dE+\tfrac12\frac{\partial E}{\partial q}\,dp
-\frac{\partial F}{\partial p}\,dp-\tfrac12\frac{\partial G}{\partial p}\,dq.
\end{equation*}

Ganz ohne Vorgänger ist \textsc{Gauss} allerdings auch hier nicht gewesen. Vielmehr
hatte \textsc{Euler}\footnote{L. \textsc{Euler}, \textit{Methodus inveniendi lineas curvas maximi minimive proprietate gaudentes}, Lausanne
1744, S. 138.} nicht nur bereits 1744 für den speziellen Fall, dass die
Fläche in der Form
\[
z=f(x,y)
\]
gegeben ist, die der Gleichung (5) entsprechende Differentialgleichung aus
dem Variationsproblem
\[
\int\sqrt{dx^2+dy^2+(T\,dx+V\,dy)^2}=\text{Minimum},
\]
wo
\[
dz=T\,dx+V\,dy
\]
gesetzt ist, abgeleitet, sondern er hatte überdies in einer späteren Arbeit\footnote{L. \textsc{Euler}, \textit{Accuratior evolutio problematis de linea brevissima in superficie quacunque ducenda},
Nova Acta Acad. Petropol., Bd. 15 (1799—1802), 1806, S. 46 (vorgelegt am 25. Jan. 1779).}
(1779) die Differentialgleichung auf die merkwürdige Form gebracht
\begin{equation*}\tag{7}
\frac{ds}{1+s^2}=\frac{dk}{(1+k^2)\sqrt{1+h^2}},
\end{equation*}
\clearpage\SourceLeaf{512}
\ArithmeticHead{55}{R}{GEODÄTISCHE POLAR- UND PARALLELKOORDINATEN.}
\noindent wo
\[
k=\frac{V}{T},\quad h^2=T^2+V^2,\quad s\sqrt{1+h^2}=\frac{V\,dx-T\,dy}{T\,dx+V\,dy}.
\]

Nun hat schon M. \textsc{Cantor}\footnote{M. \textsc{Cantor}, \textit{Vorlesungen über Geschichte der Mathematik}, Bd. 4, Leipzig 1908, S. 539. \textsc{Cantor}
schreibt $\operatorname{tg} w$ statt $\operatorname{cotg} w$.} darauf aufmerksam gemacht, dass $s=\operatorname{cotg} w$
ist, wenn $w$ den Winkel der geodätischen Linie mit der durch den betrachteten
Punkt gehenden Kurve der Schar $z=\text{const.}$ bedeutet. Und in der Tat wird
die \textsc{Gauss}sche Formel (6) für den speziellen Fall einer Parameterdarstellung,
bei der $z=q$ ist, mit der \textsc{Euler}schen Formel (7) identisch.
\begin{center}\textls[100]{B. Die Sätze über geodätische Polar- und Parallelkoordinaten.}\end{center}
\noindent\hspace*{\parindent}5. Die angeführten Stellen sind die einzigen, an denen \textsc{Gauss} selbst
unmittelbar die Variationsrechnung anwendet. Damit ist aber nur der allerkleinste
Teil der Beziehungen zwischen den \textit{Disquisitiones} und der Variationsrechnung
angegeben.

Denn da das Problem der geodätischen Linien in der Form (4) als spezieller
Fall in dem allgemeineren Variationsproblem enthalten ist, das Integral
\begin{equation*}\tag{8}
J=\int F(x,y,x',y')\,dt,
\end{equation*}
\[
x'=\frac{dx}{dt},\quad y'=\frac{dy}{dt},
\]
(wo die Funktion $F$ in Bezug auf die Variabeln $x',y'$ positiv-homogen von
der ersten Dimension ist), durch eine Kurve
\[
x=x(t),\quad y=y(t)
\]
zu einem Extremum zu machen, so lassen sich alle Sätze über geodätische
Linien in zwei Klassen einteilen:
\par\noindent\hspace*{\parindent}1) solche, die spezielle Fälle von entsprechenden allgemeinen Sätzen über
das genannte Variationsproblem sind, und
\par\noindent\hspace*{\parindent}2) solche, bei denen dies nicht der Fall ist.

Nun gehören eine Reihe von \textsc{Gauss}schen Sätzen über geodätische Linien
in die erste Kategorie, und diese haben dann befruchtend auf die Variationsrechnung
eingewirkt, insofern sie von späteren Autoren als spezielle Fälle von
\clearpage\SourceLeaf{513}
\ArithmeticHead{56}{L}{BOLZA, GAUSS UND DIE VARIATIONSRECHNUNG.}
\noindent allgemeinen Variationssätzen erkannt worden sind. Daran schliessen sich dann
weiter solche Sätze, die, ohne direkt Verallgemeinerungen von \textsc{Gauss}schen
Sätzen zu sein, als »verwandte« Sätze zu bezeichnen sind, insofern sie mit
jenen zusammen als spezielle Fälle unter ein und demselben allgemeinen
Variationssatz enthalten sind.
\par\clearpage\endgroup
